\documentclass[11pt]{article}
\usepackage[letterpaper,margin=1.0in]{geometry}
\usepackage[T1]{fontenc}
\usepackage{lmodern,amsmath,amssymb,amsthm,mathtools,booktabs,array,microtype}
\usepackage[hidelinks]{hyperref}
\usepackage{enumitem}
\usepackage{xcolor}
\definecolor{navy}{RGB}{25,50,80}
\hypersetup{colorlinks=true,linkcolor=navy,urlcolor=navy,citecolor=navy,pdftitle={Sharp torsion selection across nonunit power patterns}}
\setlength{\parskip}{3pt}
\setlength{\emergencystretch}{2em}
\setlength{\parindent}{15pt}
\newtheorem{theorem}{Theorem}[section]
\newtheorem{lemma}[theorem]{Lemma}
\newtheorem{proposition}[theorem]{Proposition}
\newtheorem{corollary}[theorem]{Corollary}
\theoremstyle{definition}
\newtheorem{definition}[theorem]{Definition}
\newtheorem{remark}[theorem]{Remark}
\newcommand{\Cp}{\mathbb C_p}
\newcommand{\Qp}{\mathbb Q_p}
\newcommand{\OO}{\mathcal O}
\newcommand{\Z}{\mathbb Z}
\newcommand{\N}{\mathbb N}
\newcommand{\ord}{\operatorname{ord}}
\newcommand{\vp}{v_p}
\newcommand{\Gauss}{G}
\newcommand{\cost}{C}
\newcommand{\supp}{\operatorname{supp}}
\newcommand{\rank}{\operatorname{rank}}
\newcommand{\ceil}[1]{\left\lceil #1\right\rceil}
\newcommand{\floor}[1]{\left\lfloor #1\right\rfloor}
\title{\textbf{Sharp torsion selection across nonunit power patterns}\\[5pt]\large Exact obstructions, finite-state covering, and dimension-free conductor asymptotics}
\author{Research note for the ProveIt project}
\date{27 November 2026}
\begin{document}
\maketitle
\begin{abstract}
The torsion-specialization argument in the proposed Dirichlet GRH development leads to a local extremal problem: how large must a valuation budget be to obstruct simultaneous nonvanishing of integral power series at a torsion tuple and at prescribed powers of that tuple? The merged ProveIt article explicitly leaves open the case of powers divisible by the residue characteristic. We give an exact weighted hitting-set characterization for arbitrary powers in one variable, a closed solution for any two powers, and an exact finite algorithm in every dimension. For a finite pattern of pure $p$-power exponents $E$, the one-variable problem reduces to covering a finite interval by translates of $-E$. More substantially, we prove that its minimum mean cycle density $\delta(E)$ is the sharp leading obstruction constant in \emph{every} dimension, with error $O_E(\sqrt N)$ for a window of $N$ torsion levels. The proof amplifies finite covering constraints into multiplicity on a smaller product grid. For $E=\{0,1,3\}$ we obtain $\delta=2/5$ and the exact one-variable cost $\lceil 2N/5\rceil$; the optimal leading conductor-exponent coefficient is consequently $5/2$, rather than the elementary $3$. For $p=2$ and the first two powers, we determine the one-variable endpoint, sharp bounds in all dimensions, and the exact endpoint in all dimensions for unramified coefficients at starting level one. All results are unconditional and concern nonvanishing, not the stronger Selmer valuation estimates. Exact arithmetic scripts accompany the note.
\end{abstract}
\tableofcontents

\section{Research context and main results}

\subsection{Which part of the repository is being extended?}

The current \emph{GeneralizedRiemannHypothesis} directory of ProveIt aggregates 209 research reports in a 982-page article \cite{project,article}. It treats the external Dirichlet RH/GRH development at revision \texttt{72fead102f663d5c7a407964a13dc0cf9d977207} as input and distinguishes unconditional method transfers from conditional analytic consequences and unresolved reductions. Its introduction identifies four transferable mechanisms: comparison of covolumes, spanning, selection of torsion twists with controlled valuations, and a one-sided analytic criterion for excluding zeros. We investigate the third.

The relevant upstream file is \texttt{New/Congruences/Selmer/RootsOfUnity.lean} \cite{roots}. It develops a degree-independent choice of roots of unity by counting roots and using a harmonic sum of torsion radii. The declarations \nolinkurl{exists_rootOfUnity_eval_le_gaussVal_add} and \nolinkurl{exists_rootsOfUnity_ellVal_eval_le_add} implement the one-variable and multivariable valuation-selection steps. The application to anticyclotomic twists is described in the merged article's source section \cite{source}. The input theorem concerns Dirichlet $L$-functions, not arbitrary nonabelian Artin or automorphic $L$-functions. We have not rebuilt or audited the upstream Lean proof; the theorems here do not depend on its correctness.

Section~V of the merged article, particularly \texttt{sections/tors.tex} \cite{tors}, establishes sharp obstruction costs for single torsion shells and for powers prime to $p$. Its theorem labeled \texttt{tors:thm:cover} gives the exact mixed-shell cost in one variable for unit powers and dimension-independent prefix bounds. Its question \texttt{tors:q:powers} asks, in particular, what happens when a prescribed power is divisible by $p$, including $p=2$ and the powers $1,2$. That is the question addressed here. The point is that a unit power permutes each shell, whereas multiplication of the exponent by $p$ descends to the parent shell, so the covering problem is no longer an independent problem on each shell.

\subsection{Summary of contributions}

\begin{enumerate}[leftmargin=*,itemsep=3pt]
\item \textbf{Exact nonunit obstruction in one variable.} For every finite set of input levels and any positive powers $J$, a finite weighted hitting problem on the torsion points gives the exact obstruction budget (Theorem~\ref{thm:one}). For two powers with different $p$-adic valuations it reduces further to a weighted vertex cover on the levels (Theorem~\ref{thm:pair}). For powers of the same valuation it is a union of cycle-cover problems. Both formulas include powers collapsing a root to the identity.
\item \textbf{Exact finite algorithm in arbitrary dimension.} Integral division reduces the optimization to a finite polynomial box. A valuation-minimal column elimination computes the least Gauss valuation in each evaluation kernel (Theorem~\ref{thm:finite}). This is an exact characterization, albeit not an efficient algorithm in large dimensions.
\item \textbf{Sharp asymptotic in every dimension.} For $J=\{p^a:a\in E\}$ the optimal leading cost is $\delta(E)N$, uniformly in dimension, where $\delta(E)$ is computable on a finite graph of $2^{\max E}$ states (Theorem~\ref{thm:main}). The error in our general multidimensional bound is $O_E(\sqrt N)$, versus $O_E(1)$ in one variable. The elementary bound $N/|E|$ is asymptotically sharp precisely when $E$ tiles the integers by translations (Corollary~\ref{cor:tile}).
\item \textbf{Explicit strict improvement.} For $E=\{0,1,3\}$ an eight-state dual certificate proves $\delta(E)=2/5$; the one-variable cost is exactly $\lceil 2N/5\rceil$ (Theorem~\ref{thm:013}). This improves the leading selection bound from $3B$ to $(5/2)B+O(\sqrt B)$ in all dimensions.
\item \textbf{The first two powers at $p=2$.} The exact one-variable endpoint is $h_0+2\lfloor B\rfloor$ in the normalization of the source question. All dimensions satisfy the bracket in Corollary~\ref{cor:two}; in particular this settles all integer budgets and all budgets with fractional part less than $1/2$. For unramified coefficients and $h_0=1$, the lower endpoint is always exact. A finite computation settles another nontrivial two-dimensional instance at $h_0=3$.
\end{enumerate}

The new uniform asymptotic does not rely on an unproved distribution of primes, any Riemann hypothesis, or the Selmer interpretation. It is a result about bounded integral series. It does not establish a new global conjecture, and the full all-dimensional exact endpoint for arbitrary ramified coefficients remains open. The comparison to prior work here is to the inspected repository; no claim of exhaustive worldwide priority is made.

\section{Local analytic setup}\label{sec:setup}

Fix a prime $p$. Write $v$ for the valuation of $\Cp$, normalized by $v(p)=1$, and put $v(0)=+\infty$ and $\OO=\{x:v(x)\ge0\}$. For $h\ge1$ define
\[
 Z_h=\{\zeta-1:\ord(\zeta)=p^h\},\qquad
 e_h=\varphi(p^h),\qquad r_h=e_h^{-1}.
\]
Every member of $Z_h$ has valuation $r_h$; $|Z_h|=e_h$. Set $Z_0=\{0\}$, the identity in the shifted coordinate. For a finite nonempty set $S\subset\Z_{\ge1}$ let $A_S=\bigcup_{h\in S}Z_h$. A positive integer $j$ acts by
\[
 [j]z=((1+z_1)^j-1,\ldots,(1+z_s)^j-1).
\]
For a nonzero integral bounded series $F=\sum_{\alpha\in\N^s}c_\alpha T^\alpha$ and depth $r>0$, its Gauss valuation is
\begin{equation}\label{eq:gauss}
 G_r(F)=\min_{\alpha}\{v(c_\alpha)+r|\alpha|\}.
\end{equation}
Set $G_r(0)=+\infty$. A witness of depth $r$ is a tuple $x$ with $v(x_i)\ge r$ for all $i$. The point $0$ is permitted. We take $r$ rational for the exact budget statements; the lower bounds also hold for real $r$.

\begin{definition}[Blocking cost]\label{def:cost}
Let $J$ be a finite nonempty set of positive integers. A finite family $F_i\in\OO[[T_1,\dots,T_s]]$, with witnesses $x_i$ of depth $r$ at which their values are nonzero, blocks $A_S^s$ at powers $J$ if
\[
 \text{for every }z\in A_S^s,\quad
 \text{there exist }i\text{ and }j\in J\text{ such that }F_i([j]z)=0.
\]
Write $C_s(S,J;r)$ for the infimum of $\sum_i v(F_i(x_i))$ over blocking families. For an interval $S=[h_0,H]\cap\Z$, the least endpoint guaranteeing simultaneous nonvanishing at budget at most $B$ is denoted $H_s(B;h_0,J,r)$.
\end{definition}

\begin{lemma}[Reduction to Gauss valuation]\label{lem:reduce}
For nonzero integral series, $G_r$ is multiplicative in the additive valuation convention, monotone in $r$, and $v(F(x))\ge G_r(F)$ at every witness of depth $r$. If $r$ is rational, some such witness attains equality. Consequently $C_s(S,J;r)$ is the infimum of $G_r(F)$ for single nonzero integral series $F$ blocking the grid.
\end{lemma}
\begin{proof}
The evaluation inequality is the ultrametric inequality, and monotonicity is immediate from the minimum formula. Choose $a,b\in\Cp^\times$ with $v(a)=r$, $v(b)=G_r(F)$. The normalized series $b^{-1}F(aY)$ has integral coefficients, tending to zero as total degree tends to infinity. Its reduction is a nonzero polynomial over the infinite residue field $\overline{\mathbb F}_p$. Choose residues of $Y$ at which it does not vanish, and lift them to integral elements. This gives the desired witness. The product of two nonzero reduction polynomials is nonzero, which proves multiplicativity. For real $r>0$ multiplicativity also follows by continuity: on any compact positive interval, the minimum in \eqref{eq:gauss} is taken among finitely many relevant terms. Now replace a family by $F=\prod_iF_i$; it has exactly the required zero-covering property and
\[
 G_r(F)=\sum_i G_r(F_i)\le\sum_i v(F_i(x_i)).
\]
Conversely a single blocking series admits a witness of depth $r$ attaining its Gauss valuation. This proves the reduction.
\end{proof}

We will also need a precise composition rule for the nonunit power.
\begin{lemma}[Gauss transport under the $p$-power map]\label{lem:transport}
Define $\theta(\rho)=\min(p\rho,1+\rho)$. For rational $\rho>0$,
\[
 G_\rho(F([p]T))=G_{\theta(\rho)}(F),\qquad
 G_\rho(F([p^a]T))=G_{\theta^{\circ a}(\rho)}(F).
\]
In particular $\theta(r_h)=r_{h-1}$ for $h\ge2$, and $\theta(r_1)=r_1+1$.
\end{lemma}
\begin{proof}
For $1\le k<p$, the coefficient of $T^k$ in $(1+T)^p-1$ has valuation one, whereas that of $T^p$ is one as an integer, of valuation zero. The Gauss valuation of this polynomial is therefore $\theta(\rho)$. Scale the variables to depth $\rho$ and the output by an element of valuation $\theta(\rho)$. The resulting normalized polynomial in one variable has a nonconstant reduction, with zero constant term. In each of the $s$ independent variables we get such a polynomial. Substitution of these nonconstant one-variable polynomials into a nonzero reduction polynomial of $F$ cannot give zero: the induced homomorphism of polynomial rings over a field is injective, as can be seen by comparing leading degrees in each variable. This proves the equality even at the radius where two terms tie. Iteration gives the second identity. The stated radius formulas follow from $e_h=(p-1)p^{h-1}$.
\end{proof}

The following grid lemma is the integral form of the multiplicity argument used in \cite{tors}. We include the proof because it is the key link from combinatorics to valuations.
\begin{lemma}[Integral product-grid multiplicity]\label{lem:grid}
Let $g\in\OO[T]$ be monic, with distinct roots forming a finite set $A$ of elements of positive valuation. Assume its nonleading coefficients have valuations bounded below by some positive number. If $P\in\OO[[T_1,\ldots,T_s]]$ vanishes to order at least $m$ at every point of $A^s$, then
\[
 P\in(g(T_1),\ldots,g(T_s))^m\OO[[T_1,\ldots,T_s]],
 \qquad G_\rho(P)\ge mG_\rho(g)\quad(\rho>0).
\]
Here vanishing to order $m$ means all Taylor terms of total degree less than $m$ are zero.
\end{lemma}
\begin{proof}
First, if $g=T^d+a(T)$ with $\deg a<d$, divide an integral series in $T_1$ by $g$. On bounded coefficient arrays let $S_d$ delete terms of $T_1$-degree less than $d$ and shift the others down by $d$. If $M_a$ denotes multiplication by $a(T_1)$, the operator $S_dM_a$ has norm strictly less than one for the coefficient supremum norm. Thus
\[
 Q=(1+S_dM_a)^{-1}S_dP=\sum_{k\ge0}(-S_dM_a)^kS_dP
\]
is integral and convergent in that complete norm, and $P-g(T_1)Q$ has degree less than $d$ in $T_1$. The same argument applies to $g^m$. Dividing successively by $g(T_i)^m$ gives integral quotients and a remainder polynomial of degree less than $md$ in every variable. Repeated monic division writes the remainder uniquely as
\[
 R=\sum_{0\le\beta_i<m} S_\beta(T)\prod_i g(T_i)^{\beta_i},\qquad
 \deg_{T_i}S_\beta<d.
\]
The quotient terms vanish to order at least $m$ on $A^s$, so $R$ does too. Induct on $k=|\beta|<m$. At a point $a\in A^s$, after lower terms have disappeared, the degree-$k$ Taylor part of $R$ is
\[
 \sum_{|\beta|=k} S_\beta(a)\prod_i\bigl(g'(a_i)(T_i-a_i)\bigr)^{\beta_i}.
\]
All these distinct monomials have zero coefficients, and $g'(a_i)\ne0$. Thus each $S_\beta$ of total index $k$ vanishes on $A^s$. The degree bounds and successive univariate interpolation force it to be the zero polynomial. Induction leaves only terms with $|\beta|\ge m$. This proves the ideal inclusion; the Gauss bound follows from integrality of all multipliers and Lemma~\ref{lem:reduce}.
\end{proof}

For $g=\prod_{h\in S}\Phi_{p^h}(1+T)$, the hypotheses hold: each cyclotomic polynomial is monic and congruent to $T^{e_h}$ modulo $p$. The product of its linear factors gives
\begin{equation}\label{eq:gridcost}
 G_\rho(g)=\sum_{h\in S}e_h\min(\rho,r_h).
\end{equation}
In particular this equals $|S|$ if $\rho\ge\max_{h\in S}r_h$.

\section{Exact characterization for arbitrary powers in one variable}

\begin{theorem}[Weighted torsion cover]\label{thm:one}
Let $s=1$. Form the finite set $V=\bigcup_{j\in J}[j]A_S$. For each $z\in A_S$ form a hyperedge
\[
 E_z=\{[j]z:j\in J\}\subseteq V.
\]
Give a vertex $a$ the weight $w(a)=\min(r,v(a))$, including $w(0)=r$. Then
\begin{equation}\label{eq:cover}
 C_1(S,J;r)=\min_{D\subseteq V\, :\, D\cap E_z\ne\varnothing\ (z\in A_S)}\ \sum_{a\in D}w(a).
\end{equation}
For rational $r$ the minimum is attained by a single polynomial.
\end{theorem}
\begin{proof}
For a blocking series $F$, its zero set $D$ in $V$ meets every edge. If $F(a)=0$ at a point of positive valuation (including $a=0$), then $F=(T-a)Q$ with $Q$ integral: the coefficients of $Q$ are convergent sums of the original coefficients times powers of $a$. Distinct zeros persist after division, so extract one linear factor for each $a\in D$. Multiplicativity gives $G_r(F)\ge\sum_{a\in D}\min(r,v(a))$. Conversely, for a hitting set $D$, the polynomial $\prod_{a\in D}(T-a)$ blocks and has Gauss valuation exactly that sum. Use Lemma~\ref{lem:reduce} to attain it with a witness. For arbitrary real $r$ the same formula holds for the infimum, by approximating $r$ from above by rational depths; attainment at the infimum need not hold in that generality.
\end{proof}

The formula is an exact finite integer optimization, not a bound involving the number or degrees of the series. One may encode all vertices by exponents modulo $p^{\max S}$ and compute their orders arithmetically, without constructing $p$-adic numbers. When all $j$ are units it specializes to the independent shell cover of \cite{tors}.

\subsection{Two arbitrary powers}
For a level $t$ let its total weight be
\begin{equation}\label{eq:weights}
 W_t=e_t\min(r,r_t)=\min(1,r/r_t)\quad(t\ge1),\qquad W_0=r.
\end{equation}

\begin{theorem}[Exact two-power formula]\label{thm:pair}
Let $J=\{j,k\}$ with $b=v_p(j)\le c=v_p(k)$. Write $j=p^b u$, $k=p^c w$, where $u,w$ are units. Put $T=\{\max(h-b,0):h\in S\}$.
\begin{enumerate}[label=(\alph*),leftmargin=*]
\item If $b=c$, let $d_t=\ord_{(\Z/p^t\Z)^\times}(wu^{-1})$ for $t\in T\setminus\{0\}$. Then
\[
 C_1(S,J;r)=\mathbf1_{0\in T}r+\sum_{t\in T\setminus\{0\}} W_t\frac{\lceil d_t/2\rceil}{d_t}.
\]
\item If $b<c$, put $a=c-b$. Construct the graph on the output levels with edges $\{t,\max(t-a,0)\}$ for $t\in T$, interpreting a singleton as a forced vertex. Then $C_1(S,J;r)$ is exactly the minimum weight of a vertex cover of this graph, with vertex weights $W_t$.
\end{enumerate}
\end{theorem}
\begin{proof}
The first power maps each input shell surjectively onto $Z_t$ for $t=\max(h-b,0)$. On $Z_t$ the second output is obtained from the first by the action of $p^{c-b}wu^{-1}\in\Z_p$ on the roots of unity, well-defined modulo the relevant order. If $b=c$, for $t>0$ this action partitions $Z_t$ into cycles of length $d_t$. Covering the edges of a cycle requires $\lceil d_t/2\rceil$ vertices (also true for a loop and for a two-cycle). Multiplying the fraction of selected vertices by the total weight $W_t$ proves (a), and level zero is a singleton.

For (b), put $\lambda=wu^{-1}\in\Z_p^\times$. Choose primitive roots $\eta_t$ for all output levels such that $\eta_t^{p^a\lambda}=\eta_{t-a}$ when $t>a$; this is possible recursively, since the map of primitive shells is surjective. For $1\le t\le a$ the image is the identity, and let $\eta_0=1$. Sample $q$ uniformly from $(\Z/p^{\max T}\Z)^\times$ if $\max T\ge1$ (if not, the claim is trivial). A vertex $\eta_t^q-1$ is uniform on $Z_t$. For each fixed $q$, a torsion hitting set $D$ selects a set of levels $L(q)=\{t:\eta_t^q-1\in D\}$ that meets every edge of the level graph. Therefore
\[
 \sum_{z\in D}w(z)=\mathbb E_q\sum_{t\in L(q)}W_t
 \ge\min_{L\text{ a level cover}}\sum_{t\in L}W_t.
\]
Selecting all torsion points of the levels in a minimum level cover gives equality. Theorem~\ref{thm:one} completes the proof.
\end{proof}

This result includes arbitrary unit factors and arbitrary input sets. The level graph is a union of chains of step $a$, possibly sharing the vertex zero; the standard two-state dynamic program on a tree computes its minimum vertex cover. For a chain of $n$ edges with every vertex of weight one the answer is $\lceil n/2\rceil$. Indeed a matching of that size gives a lower bound and alternating vertices give equality. For an interval of positive levels with no collapse and all weights one, the chains are the residue classes modulo $a$, so the cost is the sum of $\lceil n_i/2\rceil$ over their edge counts $n_i$.

\subsection{All pure $p$-power patterns}
\begin{theorem}[Level hypergraph]\label{thm:levels}
Let $J=\{p^a:a\in E\}$ for a finite nonempty $E\subseteq\Z_{\ge0}$. For every $S$, $C_1(S,J;r)$ is the minimum weight of a set of levels meeting every hyperedge
\[
 \{\max(h-a,0):a\in E\}\qquad(h\in S),
\]
with weights \eqref{eq:weights}. The same formula holds if all powers are multiplied by the same integer prime to $p$.
\end{theorem}
\begin{proof}
Take compatible primitive roots $\eta_t$ with $\eta_t^p=\eta_{t-1}$ and $\eta_0=1$. Sample a common uniform unit $q$ as in Theorem~\ref{thm:pair}. Along the chain $\eta_h^q$ all power outputs belong to the same coherent branch; a torsion hitting set selects levels covering every edge. Averaging the weights proves the lower bound, and taking entire shells proves equality. Multiplying all powers by the same unit just permutes the input tuples.
\end{proof}

\section{An exact finite computation in arbitrary dimension}\label{sec:finite}

The reduction in this section applies equally to the mixed-shell unit-power question and to the new nonunit-power question. It gives an effective answer for any specified finite instance, even though a concise general closed form may not be available.

Let $A'=\bigcup_{j\in J}[j]A_S$ in one variable, let $g(T)=\prod_{a\in A'}(T-a)$, let $d=\deg g$, and define the box of monomials
\[
 \mathcal M=\{T^\alpha:0\le\alpha_i<d\},\qquad m_0=d^s.
\]
Every output tuple lies in $(A')^s$. Write $V_s=\bigcup_{j\in J}[j]A_S^s$ and form edges $\{[j]z:j\in J\}$ for $z\in A_S^s$. For a hitting set $D\subseteq V_s$ form the evaluation matrix $A_D=(a^\alpha)_{a\in D,\alpha\in\mathcal M}$ over the cyclotomic local field $K=\Qp(\mu_{p^{\max S}})$. Take an integral basis $P_1,\dots,P_k$ of its coefficient kernel, where integral means a basis over the valuation ring $\OO_K$ for the kernel in $\OO_K^{m_0}$, not an orthogonal basis for the weighted Gauss norm.

\begin{theorem}[Finite exact obstruction algorithm]\label{thm:finite}
With these definitions, for rational $r>0$,
\begin{equation}\label{eq:finite}
 C_s(S,J;r)=\min\left\{G_r(g),\ \min_{D\text{ hitting}}\min_{1\le i\le k(D)}G_r(P_i)\right\},
\end{equation}
where an empty kernel contributes $+\infty$. It suffices to enumerate inclusion-minimal hitting sets. The minimum is attained by one integral polynomial over $K$ with a witness in $\Cp^s$. All the linear algebra can be done exactly in the rational cyclotomic number field embedded in $K$.
\end{theorem}
\begin{proof}
The set $A'$ is a union of full primitive shells and possibly $\{0\}$. Hence $g$ is monic and congruent to $T^d$ modulo $p$. Integral division as in Lemma~\ref{lem:grid} gives
\[
 F=\sum_i g(T_i)Q_i+R,\qquad R\in\operatorname{span}_{\OO}\mathcal M.
\]
If $G_r(F)<G_r(g)$ then $G_r(R)=G_r(F)$, and $R$ vanishes wherever $F$ vanishes on $V_s$. Thus every blocker with smaller Gauss cost than $g$ has a polynomial remainder in one of the hitting-set kernels. If its cost is not smaller, $g(T_1)$ itself is a blocker of cost $G_r(g)$.

It remains to justify the kernel calculation even for coefficients in $\OO=\OO_{\Cp}$. Over the discrete valuation ring $\OO_K$, invertible row and column operations put $A_D$ in Smith form with nonzero diagonal entries followed by zeros. The same operations are invertible over $\OO$. Since $\OO$ is a domain, the coefficients in the nonzero diagonal positions must be zero, and every integral kernel vector over $\OO$ is an integral combination of the remaining basis vectors. Hence for its polynomial $P=\sum_i b_iP_i$,
\[
 G_r(P)\ge\min_i\{v(b_i)+G_r(P_i)\}\ge\min_i G_r(P_i).
\]
The latter minimum is achieved by an individual basis vector. Lemma~\ref{lem:reduce} supplies the witness, and every polynomial considered on the right of \eqref{eq:finite} blocks. Any hitting set contains a minimal one, so the restricted enumeration suffices. Matrix entries and elimination operations are exact in $\mathbb Q(\zeta_{p^{\max S}})$, using the valuation of its chosen embedding, because a nonzero minimum-valuation pivot divides all entries to be eliminated in the local valuation ring. This proves the formula and effective attainment.
\end{proof}

A practical kernel algorithm requires only column elimination. Start with columns representing the monomials and their evaluation vectors. Choose a nonzero entry of least valuation in the remaining evaluation matrix. Subtract integral multiples of its column from every other column to clear that row; then discard the pivot row and pivot column. Every discarded coefficient is forced to zero in the transformed kernel. Repeat. When the remaining evaluation vectors are zero, the surviving polynomial columns form the desired integral basis. The column operations are invertible, so no solution is lost. For $K=\Qp(\zeta_{p^H})$, $\pi=\zeta_{p^H}-1$ is a uniformizer and the valuations of the terms $q_i\pi^i$, $0\le i<e_H$, are distinct modulo one when nonzero. Thus
\[
 v\left(\sum_{i=0}^{e_H-1}q_i\pi^i\right)
 =\min_i\{v_p(q_i)+i/e_H\}\qquad(q_i\in\mathbb Q).
\]
This gives an exact pivot oracle. The procedure is exponential or worse in the input grid size, as expected for enumeration of hitting sets; its usefulness is certification of small cases and possible testing of closed-form conjectures.

\section{Covering words and their sharp asymptotic density}\label{sec:words}

We now focus on a finite set $E\subseteq\Z_{\ge0}$ with $0\in E$, and put $L=\max E$, $k=|E|$. Define
\begin{equation}\label{eq:D}
 D_M(E)=\min\left\{\sum_{i=0}^{M-1} b_i:\ b_i\in\{0,1\},\ \sum_{a\in E} b_{t+a}\ge1\ (0\le t<M-L)\right\}
\end{equation}
for $M\ge L+1$. This is the minimum number of marked positions needed to hit every translate of $E$ wholly inside a word of length $M$. For $E=\{0\}$ set $D_M=M$.

Consider the directed graph with states $u\in\{0,1\}^L$. An edge appends one bit $d$, discards the oldest bit, and has weight $d$; it is allowed exactly when the resulting $(L+1)$-bit block hits $E$. For $L=0$ use a single vertex with a weight-one loop. For $L>0$ the graph is strongly connected: append $L$ ones to reach the all-ones state, then append the desired state, which is permitted because the first bit of each intervening block is one. Define $\delta(E)$ as the minimum mean weight of a directed cycle in this graph.

\begin{proposition}[Finite-state covering constant]\label{prop:mean}
For $L\ge1$, let $K_E=L+2^L-1$. Then
\begin{equation}\label{eq:wordbound}
 \delta(E)(M-K_E)\le D_M(E)\le\delta(E)M,
 \qquad \frac1k\le\delta(E)\le1.
\end{equation}
The value $\delta(E)$ is rational and computable by a finite graph algorithm. An optimal infinite periodic marking exists.
\end{proposition}
\begin{proof}
A minimum mean cycle is attained by a simple cycle, since any nonsimple cycle decomposes into shorter cycles and its mean is their length-weighted mean. Repeating it produces a periodic binary sequence, every $E$-translate of which is hit. Average the number of ones in a word of length $M$ over all shifts of one period; at least one word has at most $\delta M$ ones. For the reverse inequality, the allowed word gives a path of length $M-L$. Remove cycles until at most $2^L-1$ edges remain. All removed cycles have mean at least $\delta$, and all remaining weights and initial bits are nonnegative. Thus the path weight is at least $\delta(M-L-2^L+1)$, proving the lower bound. Averaging the constraints of a periodic sequence over one period shows $k\delta\ge1$. The all-ones word is feasible. Rationality and effective computation follow from the finite graph (for example, by enumerating its simple cycles or using a minimum-mean-cycle algorithm).
\end{proof}

\begin{corollary}[When the elementary density is sharp]\label{cor:tile}
The equality $\delta(E)=1/k$ holds if and only if $E$ tiles $\Z$ by translations.
\end{corollary}
\begin{proof}
If equality holds, the periodic minimum cycle gives $\sum_{a\in E}b_{t+a}\ge1$ with average one, hence equality at every $t$. For the set $D$ of marked positions, the sets $d-E$ ($d\in D$) are pairwise disjoint and partition $\Z$; reflection gives a tiling by $E$. Conversely, a tiling by $-E$ gives a binary marking with exactly one hit at every translate. Counting over any long interval of length $M$ gives $M/k+O_E(1)$ marked positions, because only $O_E(1)$ terms at the ends differ when summing the shifted counts. Applying the lower bound of \eqref{eq:wordbound} and letting $M\to\infty$ yields $\delta\le1/k$. The opposite inequality was proved above.
\end{proof}

\section{Dimension-free sharp leading cost}\label{sec:main}

Let $S=[h_0,H]\cap\Z$ with $N=H-h_0+1\ge1$, let $J=\{p^a:a\in E\}$, and assume
\begin{equation}\label{eq:deep}
 r\ge\theta^{\circ L}(r_{h_0}).
\end{equation}
When $h_0>L$, this is just $r\ge r_{h_0-L}$: every nontrivial output shell is within the allowed witness depth. The more general condition also handles collapsed outputs. It is not asserted that these depth conditions are necessary for the bounds below.

\begin{theorem}[Chain amplification]\label{thm:main}
Under \eqref{eq:deep}, for every $s\ge1$,
\begin{equation}\label{eq:chain}
 C_s(S,J;r)\ge \max\left\{\frac N k,\ \max_{L+1\le M\le N+L}\frac{D_M(E)}{M}(N-M+L+1)\right\}.
\end{equation}
If $h_0>L$, the exact one-variable cost and an upper bound in all dimensions are
\begin{equation}\label{eq:upper}
 C_1(S,J;r)=D_{N+L}(E),\qquad C_s(S,J;r)\le D_{N+L}(E).
\end{equation}
For arbitrary $h_0\ge1$ satisfying \eqref{eq:deep}, one has the upper bound $C_s(S,J;r)\le\delta(E)(N+L)+L$. In particular,
\begin{equation}\label{eq:asymp}
 \boxed{\displaystyle C_s([h_0,H],J;r)=\delta(E)N+O_E(\sqrt{N+1})}
\end{equation}
uniformly in $s,p,h_0,r$ satisfying the hypotheses. For $h_0>L$ the error in one variable is $O_E(1)$.
\end{theorem}
\begin{proof}
Let $F$ be a single blocking integral series of Gauss budget at most $B$. For each $z\in A_S^s$, at least one factor of
\[
 P_0(T)=\prod_{a\in E} F([p^a]T)
\]
vanishes. At $\rho=r_{h_0}$, Lemma~\ref{lem:transport} and \eqref{eq:deep} bound $G_\rho(P_0)$ above by $k B$. Lemma~\ref{lem:grid} with \eqref{eq:gridcost} bounds it below by $N$. This proves the first lower bound.

For the amplified bound, choose $M\in[L+1,N+L]$ and put $h_*=h_0+M-L-1$, $S_*=[h_*,H]\cap\Z$. It contains $N-M+L+1$ levels. For any tuple $z\in A_{S_*}^s$ and each $0\le t<M-L$, the tuple $[p^t]z$ lies in $A_S^s$: every coordinate order exponent is reduced by $t$, remains at least $h_0$, and is at most $H$. Since $F$ blocks $A_S^s$, the bits
\[
 b_j(z)=\mathbf1_{F([p^j]z)=0}\qquad (0\le j<M)
\]
satisfy all the constraints in \eqref{eq:D}. Therefore at least $D_M(E)$ of the factors of
\[
 P_M(T)=\prod_{j=0}^{M-1}F([p^j]T)
\]
vanish at $z$. Each such factor is analytic near $z$ and has zero constant Taylor term there. The product has vanishing order at least $D_M(E)$ on the entire grid $A_{S_*}^s$.

Set $\rho=r_{h_*}$. Since $h_*-h_0=M-L-1$, we have $\theta^{\circ(M-1)}(\rho)=\theta^{\circ L}(r_{h_0})\le r$. The composition lemma gives $G_\rho(P_M)\le MB$. On the other hand, Lemma~\ref{lem:grid} gives
\[
 G_\rho(P_M)\ge D_M(E)(N-M+L+1).
\]
This proves the second lower bound in \eqref{eq:chain}. This argument works with mixed coordinate levels and does not assume that all coordinates lie on the same shell.

If $h_0>L$, all output levels are positive and of total weight one by \eqref{eq:deep}. Index them in descending order $H,H-1,\dots,h_0-L$. An input level $h$ corresponds to $t=H-h$, and its outputs correspond to $t+E$. Theorem~\ref{thm:levels} gives precisely $D_{N+L}(E)$. A minimizing set of levels $D$ gives the blocker $\prod_{h\in D}\Phi_{p^h}(1+T_1)$ in every dimension, of budget $|D|$ at witness zero.

If $h_0\le L$, take a periodic minimizing marking on the same integer interval of length $N+L$, allowing its lower indices to represent nonpositive levels. Choose a shift with at most $\delta(N+L)$ marks. Discard all nonpositive levels, and add every input level $h\le L$. All other inputs have only positive outputs and hence remain covered. This uses at most $L$ additional marked levels, each of cost one because of \eqref{eq:deep} and the monotonicity of $r_h$. The resulting one-coordinate polynomial blocks as before.

For the asymptotic when $L\ge1$, choose $M$ of order $\sqrt{(K_E+1)N}$, rounded to an integer in $[L+1,N+L]$. For sufficiently large $N$ this lies in the required interval without clipping. Proposition~\ref{prop:mean} implies
\[
 \frac{D_M(E)}M(N-M+L+1)
 \ge\delta(E)N-\delta(E)\left(M+\frac{K_E N}{M}\right),
\]
by expanding the product (or dropping the positive cross term). This proves a lower bound $\delta N-O_E(\sqrt{N+1})$; finitely many smaller $N$ are covered by enlarging the same $E$-dependent constant. The upper bound has error $O_E(1)$. If $L=0$, the product-grid lemma and the one-coordinate full-shell blocker give $C_s=N$ exactly. Finally \eqref{eq:wordbound} gives the one-variable error $O_E(1)$ when $h_0>L$.
\end{proof}

\begin{corollary}[Conductor-exponent law]\label{cor:endpoint}
For fixed $E$ and the depth condition \eqref{eq:deep}, as $B\to\infty$,
\begin{equation}\label{eq:endpoint}
 H_s(B;h_0,J,r)=h_0+\frac B{\delta(E)}+O_E(\sqrt{B+1}),
\end{equation}
uniformly in $s,p,h_0,r$. For $s=1$ and $h_0>L$ the error is $O_E(1)$. The elementary leading upper coefficient $k=|E|$ is optimal if and only if $E$ tiles $\Z$.
\end{corollary}
\begin{proof}
The cost is nondecreasing in the window length. Theorem~\ref{thm:finite} shows that for rational $r$ a blocker attains the threshold cost, so the least successful window has $C_s>B$. The lower estimate in \eqref{eq:asymp} ensures success at $N=B/\delta+O_E(\sqrt{B+1})$ with a sufficiently large positive error. The explicit blockers with cost $\le\delta N+O_E(1)$ show failure up to $N=B/\delta-O_E(1)$. Since $H=h_0+N-1$, \eqref{eq:endpoint} follows. The exact one-variable upper and lower estimates give the improved error. The characterization of the coefficient follows from Corollary~\ref{cor:tile}.
\end{proof}

\begin{remark}
The new point is the use of a \emph{long chain} of constraints, not just the elementary product over the specified powers. In the chain product, a zero at one shift may cover several constraints, but the finite covering number enforces its correct reuse rate. Vanishing multiplicities then transmit the full covering density to a product grid. This is why the leading constant is independent of dimension despite the mixed-level coupling.
\end{remark}

\section{Two explicit families}
\subsection{The first two powers, including $p=2$}
\label{sec:two}
Take $E=\{0,1\}$, so $J=\{1,p\}$, with $r\ge\theta(r_{h_0})$. If $h_0\ge2$, the level graph for the interval has $N$ edges and $N+1$ vertices, all of weight one. If $h_0=1$, the bottom vertex has weight $r\ge1$, all others have weight one, and a minimum alternating cover can avoid the bottom vertex. Thus Theorem~\ref{thm:pair} gives $C_1=\lceil N/2\rceil$ in both cases. The elementary product bound already gives $C_s\ge N/2$ for every $s$.

\begin{corollary}[Exact and near-exact endpoints]\label{cor:two}
With $J=\{1,p\}$ and the depth condition above,
\[
 C_1=\ceil{N/2},\qquad \frac N2\le C_s\le\ceil{N/2},
\]
and for all $B\ge0$,
\begin{equation}\label{eq:twoendpoint}
 H_1=h_0+2\floor B,\qquad
 h_0+2\floor B\le H_s\le h_0+\floor{2B}.
\end{equation}
Hence every even window has exact cost $N/2$ in all dimensions. The two endpoint bounds coincide for $B-\lfloor B\rfloor<1/2$, and for the first two powers at $p=2$ this includes the normalization $h_0=h_*(r/2)$ of the open question in \cite{tors}, where $h_*(x)=\min\{h\ge1:r_h\le x\}$.
\end{corollary}
\begin{proof}
The costs follow from the preceding argument and the explicit one-coordinate blocker. The least $N$ with $\lceil N/2\rceil>B$ is $2\lfloor B\rfloor+1$, which gives $H_1$. In every dimension, a window of $2\lfloor B\rfloor$ levels (when nonempty) can be blocked with budget at most $B$; a window of $\lfloor2B\rfloor+1$ levels cannot, by the lower bound. For $p=2$, $r\ge 2r_{h_0}$ is exactly sufficient for $r\ge\theta(r_{h_0})$, also at $h_0=1$.
\end{proof}

\begin{corollary}[Unramified coefficients at $p=2$]\label{cor:unram}
Suppose $p=2$, $h_0=1$, $r\ge2$, and every $F_i$ has coefficients in the valuation ring of an unramified extension of $\mathbb Q_2$ with integer value group. Then the optimal endpoint in \emph{all} dimensions is $1+2\lfloor B\rfloor$.
\end{corollary}
\begin{proof}
For the product $F$, both $G_1(F)$ and $G_2(F)$ are integers and at most $B$, since all nonzero coefficient valuations and both depths are integers. If it blocks, Lemmas~\ref{lem:transport} and \ref{lem:grid} at $\rho=1$ give $N\le G_1(F)+G_2(F)\le2\lfloor B\rfloor$. The alternating shell blocker is in $\mathbb Z_2[T_1]$ and attains the lower endpoint. No assumption is made that the witnesses themselves lie in the coefficient field.
\end{proof}

For $E=\{0,1,\ldots,L\}$ with $h_0>L$, the same argument gives $D_{N+L}=\lceil N/(L+1)\rceil$. Disjoint length-$(L+1)$ constraints give the lower bound and a single mark every $L+1$ positions gives the upper bound. In particular all windows with $N$ divisible by $L+1$ have exactly that cost in every dimension. The corresponding endpoint bracket is
\[
 h_0+(L+1)\lfloor B\rfloor\le H_s\le h_0+\lfloor(L+1)B\rfloor.
\]
Thus the exact finite problem is not confined to two powers.

\subsection{A non-tiling pattern: $E=\{0,1,3\}$}\label{sec:013}

\begin{theorem}[An explicit strict improvement]\label{thm:013}
For $E=\{0,1,3\}$ one has
\[
 \delta(E)=\frac25,\qquad D_{N+3}(E)=\ceil{\frac{2N}{5}}\quad(N\ge1).
\]
For $h_0>3$, $r\ge r_{h_0-3}$ and $J=\{1,p,p^3\}$,
\[
 H_1=h_0+\floor{\tfrac52\floor B},\qquad
 H_s=h_0+\tfrac52 B+O(\sqrt{B+1})\quad(s\ge1),
\]
with an absolute implied constant independent of the prime and dimension.
\end{theorem}
\begin{proof}
The states are the three-bit strings $abc$; the allowed edge $abc\to bcd$ has cost $d$ and condition $a+b+d\ge1$. Consider the following integer potential:
\begin{center}
\begin{tabular}{c|rrrrrrrr}
state&000&001&010&011&100&101&110&111\\ \midrule
$\phi$&$-6$&$-3$&$-2$&$0$&$-4$&$0$&$-2$&$0$
\end{tabular}
\end{center}
For each of the 14 allowed edges, direct substitution gives
\begin{equation}\label{eq:dual}
 5d-2\ge\phi(bcd)-\phi(abc).
\end{equation}
Summing over a cycle shows that its mean is at least $2/5$. The cycle
\[
 000\to001\to011\to110\to100\to000
\]
has appended bits $11000$ and mean $2/5$, proving the first assertion. In particular, the pattern does not tile $\Z$, consistently with Corollary~\ref{cor:tile}.

For a word of length $N+3$, the path has $N$ edges. If $u$ is its initial state, $v$ its final state and $\ell(u)$ the number of initial ones, summing \eqref{eq:dual} gives
\[
 5\sum_{i=0}^{N+2} b_i\ge 2N+5\ell(u)+\phi(v)-\phi(u)\ge2N.
\]
The last inequality follows from $\min\phi=-6$ and $5\ell(u)-\phi(u)\ge6$ for all eight states. Hence $D_{N+3}\ge\lceil2N/5\rceil$.

For equality, write $N=5q+t$, $0\le t<5$. If $t\in\{0,1,3,4\}$, start at $000$, follow $q$ full copies of the displayed cycle, then its first $t$ edges. The number of ones equals $2q+\lceil2t/5\rceil$. If $t=2$, start at $010$ and take $010\to100\to000$, with two appended zeros and one initial one, then follow $q$ full cycles. This has $2q+1$ ones. Thus the lower bound is attained for every $N$. The endpoint formula in one variable follows from $\lceil2N/5\rceil\le B$ if and only if $N\le\lfloor(5/2)\lfloor B\rfloor\rfloor$, and the other endpoint formula follows from Corollary~\ref{cor:endpoint}.
\end{proof}

This example uses three selected powers. Counting them separately gives a sufficient endpoint coefficient of $3$ from the first bound of \eqref{eq:chain}, but the sharp asymptotic coefficient is $5/2$. The improvement is $1/6$ of the elementary coefficient; equivalently the leading obstruction cost rises from $N/3$ to $2N/5$. This is a genuine asymptotic gain over the direct union/product argument even for arbitrarily many coordinates. The exponent is logarithmic in the order $p^H$, so the change in the allowed order is exponential in $B$ at fixed $p$.

\begin{table}[htbp]
\centering
\caption{Exact one-variable cost and proved lower bounds for all dimensions for $E=\{0,1,3\}$. The depth assumption and $h_0>3$ are in force. The elementary bound is $N/3$; the amplified column is the maximum in \eqref{eq:chain}.}
\label{tab:bounds}
\begin{tabular}{rrrr}
\toprule
$N$&Elementary lower bound&Amplified lower bound&$C_1$ (upper bound for $C_s$)\\\midrule
10&$10/3$&$10/3$&4\\
20&$20/3$&$20/3$&8\\
50&$50/3$&$50/3$&20\\
100&$100/3$&$372/11$&40\\
200&$200/3$&$141/2$&80\\
1000&$1000/3$&$11676/31$&400\\
\bottomrule
\end{tabular}
\end{table}

\section{Finite checks and a small two-dimensional instance}

The accompanying \texttt{check\_patterns.py} implements the finite-word dynamic program and an exhaustive torsion-node hitting-set solver. It uses exact rational arithmetic. The tests compare the finite-state recurrence with exhaustive binary words, verify the potential \eqref{eq:dual} on every allowed edge, check the closed formula in Theorem~\ref{thm:013}, and compare level covers against direct torsion-node covers, including arbitrary pairs of powers with unit factors. The output \texttt{checks.txt} records the results. These checks are supplementary: the general theorems above have proofs not relying on finite experiments.

For a separate example of Theorem~\ref{thm:finite}, take $p=2$, $s=2$, $S=\{3\}$, $J=\{1,2\}$ and $r=1/2$. Each input tuple of exact order-eight coordinates maps by squaring to a parent tuple of exact order-four coordinates. There are four parent tuples, and each has four children. The hypergraph is the disjoint union of four stars, with every edge consisting of a child and its parent. An inclusion-minimal cover chooses either the parent alone or all four children of each star. There are exactly $2^4=16$ such covers.

The output coordinate set has six points (four of order eight and two of order four), so $d=6$ and the remainder box contains 36 monomials. The field is $K=\mathbb Q_2(\zeta_8)$, computed in the script as $\mathbb Q[X]/(X^4+1)$ with the valuation obtained after expansion in the uniformizer $X-1$. The finite algorithm gives a minimum budget of $1$; all sixteen kernels have minimum at least one. The upper bound is independently witnessed by $\Phi_4(1+T_1)$, since every squared input has first coordinate of order four. In particular $C_2(\{3\},\{1,2\};1/2)=1$. Here the claimed exact lower bound is a finite computation, whose source and output are included as \texttt{check\_small\_grid.py} and \texttt{grid\_checks.txt}. There is no appeal to approximate $p$-adic numerics.

\begin{table}[htbp]
\centering
\caption{Kernel minima for the four-star example, grouped by the number of parents selected. The individual 16 values are recorded in \texttt{grid\_checks.txt}.}
\begin{tabular}{rrl}
\toprule
Parents selected&Number of covers&Minimum Gauss valuations of their kernels\\\midrule
0&1&$1$\\
1&4&$3/2$ (all four)\\
2&6&$1$ (all six)\\
3&4&$1$ (all four)\\
4&1&$1$\\
\bottomrule
\end{tabular}
\end{table}

\section{Consequences for the twisting and determinant interface}

A direct application of the selection theorem is simultaneous generic-rank preservation. Let $A_i(T)$ be finitely many integral matrices, with a chosen nonzero minor $F_i$ of maximal generic rank in each. Suppose there are witnesses of depth $r$ with total valuation budget $B$. Choose $h_0$ satisfying \eqref{eq:deep}. For any $N$ for which the right side of \eqref{eq:chain} exceeds $B$, the family is not a blocker. Thus there is a tuple $z$ with every coordinate in the stated window such that
\[
 \rank_{\Cp} A_i([p^a]z)=\rank_{\operatorname{Frac}(\OO[[T]])} A_i(T)
 \qquad(i,\ a\in E).
\]
The selected minors are nonzero, so the ranks are at least the generic ranks, and larger minors vanish identically. In particular, a generically invertible matrix remains invertible at all the selected powers. For a finite complex of free modules one can choose maximal minors of its differentials; the resulting complex over $\Cp$ has the generic cohomology dimensions at every chosen power. This uses exactly the rank interface of \cite{tors} and improves the leading selection exponent for non-tiling pure-power patterns.

In the anticyclotomic construction of \cite{roots,source}, raising a character to a power divisible by $p$ decreases its $p$-power order. To preserve nontriviality of each coordinate character, one should take $h_0>L$; then all selected powers in this note remain nontrivial in every coordinate. The order cutoff $p^H$ is uniform in the degrees and the number of variables. To turn a rank statement into quantitative bounds on specialized Selmer lengths, however, one needs valuations of the minors, not merely nonzero values. The present theorem does not control how close $F_i([p^a]z)$ is to zero, does not identify an arithmetic Selmer matrix, and does not bound the norms of the auxiliary primes. It is not a substitute for the positive-tolerance selector used in the full source proof.

\section{Open questions and limitations}

The main unresolved finite-dimensional issue is whether the exact one-coordinate upper bound $D_{N+L}(E)$ is attained as the optimum in every dimension under the noncollapse depth hypothesis. Theorem~\ref{thm:main} proves the same leading constant, not this equality. In particular, at $p=2$ with the first two powers and arbitrary ramified coefficients, the odd-window cost remains between $N/2$ and $\lceil N/2\rceil$ in general. The finite algorithm settles any given finite case, but grows rapidly. Specific next directions are:
\begin{enumerate}[leftmargin=*,itemsep=3pt]
\item Replace the $O_E(\sqrt N)$ loss by $O_E(1)$, or determine whether a dimension-dependent negative secondary term is unavoidable. The current loss is from discarding initial shells in order to amplify the chain length. A weighted grid-multiplicity argument on nested shells might avoid this loss.
\item Determine whether multi-coordinate blockers can ever be strictly cheaper than the best one-coordinate blocker for $J=\{1,p\}$, or more generally for pure-power patterns. The four-star check and the exact even-window cases are evidence only for their respective settings.
\item For general $J$ with both unit factors and distinct $p$-adic valuations, the exact one-variable hypergraph is known by Theorem~\ref{thm:one}. Can a suitable finite-state model, perhaps with a profinite unit action, give an explicit leading constant uniform in dimension? The present chain proof needs pure $p$-power patterns.
\item Resolve the exact mixed-shell unit-power formula asked in \texttt{tors:q:powers}. Theorem~\ref{thm:finite} provides a finite algorithm but not a closed formula for that case.
\item Seek analogues for Lubin--Tate formal groups and for the stronger requirement of small valuation excess. The latter is the substantive arithmetic input needed to improve the denominator bounds in the proposed GRH proof; a proof of nonvanishing alone does not settle it.
\end{enumerate}

\paragraph{Conclusion.} We obtain a sharp dimension-free leading obstruction law by combining the upstream torsion-selection mechanism with finite-state covering and multiplicity on product grids. The new exact one-variable formulas and the explicit non-tiling example show that nonunit powers can have a fundamentally different optimal exponent from the naive product bound. This resolves concrete parts of a question explicitly posed in the repository and supplies a testable framework for the remaining exact finite problem. It does not claim resolution of a major external conjecture such as general automorphic GRH, Artin's primitive-root conjecture, or a Selmer main conjecture.

\begingroup\small
\begin{thebibliography}{9}
\setlength{\itemsep}{1pt}
\bibitem{project} ProveIt, \emph{Generalized Riemann Hypothesis}, project README. \href{https://github.com/VladimirReshetnikov/ProveIt/blob/main/Analysis/GeneralizedRiemannHypothesis/README.md}{GitHub source}. Blob \texttt{9aa123f3c17f3e8056d64d3c2bc7823e583fa82d}.
\bibitem{article} ProveIt, \emph{Beyond the Riemann hypothesis for Dirichlet $L$-functions: consequences, transfers and limits}, merged article, README. \href{https://github.com/VladimirReshetnikov/ProveIt/blob/main/Analysis/GeneralizedRiemannHypothesis/Article/README.md}{GitHub source}. Blob \texttt{b986fb54c6e641e9241439476bb5579e234fecd6}.
\bibitem{source} ProveIt, merged article, \emph{The source theorem and its transferable mechanism}, \texttt{Article/sections/src.tex}, in particular \texttt{src:sec:mechanism}. \href{https://github.com/VladimirReshetnikov/ProveIt/blob/main/Analysis/GeneralizedRiemannHypothesis/Article/sections/src.tex}{GitHub source}. Blob \texttt{b0ac2147291b7da3a74cdf6072cc557b1e1f5ee8}.
\bibitem{tors} ProveIt, merged article, \emph{Torsion specialization: selection theorems}, \texttt{Article/sections/tors.tex}. See especially labels \texttt{tors:lem:grid}, \texttt{tors:thm:cover}, \texttt{tors:prop:powers}, \texttt{tors:thm:twopowers}, and \texttt{tors:q:powers}. \href{https://github.com/VladimirReshetnikov/ProveIt/blob/main/Analysis/GeneralizedRiemannHypothesis/Article/sections/tors.tex}{GitHub source}. Blob \texttt{73112f46bf45e0d391a0642b4a8241dd340acada}.
\bibitem{roots} L. Sernau et al., \emph{grh-lean}, readable edition, revision \texttt{72fead102f663d5c7a407964a13dc0cf9d977207}, \nolinkurl{grh/src/New/Congruences/Selmer/RootsOfUnity.lean}. \href{https://github.com/lsernau/grh-lean/blob/72fead102f663d5c7a407964a13dc0cf9d977207/grh/src/New/Congruences/Selmer/RootsOfUnity.lean}{Upstream source}; \href{https://github.com/VladimirReshetnikov/ProveIt/blob/main/Analysis/GeneralizedRiemannHypothesis/vendor/grh-lean/grh/src/New/Congruences/Selmer/RootsOfUnity.lean}{vendored source}. Vendored blob \nolinkurl{b5ad128a2c44538e8783050a371e1e44af0c32d1}.
\end{thebibliography}
\endgroup
\end{document}
